\chapter{پژوهش‌های پیشین}\label{Chap:Chap2}
\minitoc
\setcounter{footnote}{0}
\pagestyle{plain}
در این فصل ابتدا به مرور کارهای مهم در زمینه یادگیری مستمر پرداخته، سپس پژوهش‌های حائز اهمیت در زمینه متا-یادگیری را معرفی کرده و سپس به کارهای پیشین در زمینه یادگیری متا-مستمر می‌پردازیم.

\section{یادگیری نمونه محدود با داده افزایی}

به طور کلی، داده‌افزایی در دو سطح بردار ویژگی‌ها و در سطح نمونه‌‌ها تعریف می‌شود. هدف داده‌افزایی تولید و گسترش مجموعه‌ی دادگان است به طوری که از توزیع کلی داده‌های آموزش خارج نشویم. چالش‌هایی مثل یافتن تبدیل‌های مناسب برای تولید نمونه‌های معنادار و داشتن سوگیری نمونه‌های تولید شده روی مجموعه‌ی داده‌های آموزش اولیه، از مشکلات این روش است.

\subsection{داده‌افزایی در سطح بردار ویژگی}

در این دسته، سعی می‌شود تبدیل‌هایی روی بازنمایی‌های بدست آمده از یک شبکه‌ی عمیق اعمال شود تا به بازنمایی داده‌هایی معنادار از همان دسته برسیم. یکی از روش‌های مورد استفاده، داده‌افزایی به وسیله‌ی
\trans{خیال‌بافی}{Hallucination}
 است
 \cite{Hariharan2016LowShotVR}
 .
 در این روش، تبدیل‌های انتزاعی که حاصل تغییرات بین نمونه‌های متعلق به یک دسته‌ی واحد است، از روی بردار ویژگی‌های آن‌ها محاسبه می‌شود. این تغییرات بدست آمده، ممکن است منحصر به یک دسته‌ی خاص باشد. برای مقابله با این موضوع، بین تمامی تغییرات درون دسته‌ای، یک معیار شباهت محاسبه می‌شود و متداول‌ترین آن‌ها انتخاب می‌شود. 
 
 به عنوان مثال، اگر در فضای ویژگی‌ها، بازنمایی دو نمونه مربوط به دسته‌ی اول را داشته باشیم و تغییرات بین آن‌ها را به صورت
 $z_1^a - z_2^a$
 محاسبه کنیم، اگر همین تغییرات را برای دو نمونه از یک دسته‌ی دیگر به صورت 
  $z_1^b- z_2^b$
  محاسبه کنیم، در صورتی این تفاضل به تبدیل‌های قابل اعمال اضافه می‌شود که شباهت کسینونسی این دو تفاضل بالا باشد. در نتیجه تبدیل‌هایی ذخیره می‌شوند که مستقل از دسته باشند. ویژگی کلی داده‌افزایی در سطح ویژگی، داشتن امکان تبدیلات پیچیده‌تر نسبت به فضای تصویر است. به عنوان مثال، در این حالت ممکن است تبدیلی را اعمال کنیم که بازنمایی یک پرنده‌ی ساکن روی درخت را به بازنمایی یک پرنده در حال پرواز مبدل سازد.
  
  \begin{figure}[!htbp]
  	\begin{center}
  		\includegraphics[width=1\textwidth]{images/imaginary1.png}
  		%		\vspace{0.5in}
  		\caption{
  			قسمت خیال‌باف با G
  			و دسته‌بند با h نشان داده شده است. برای یک اپیزود متایادگیری، به نمونه‌های پشتیبان هر دسته تعدادی داده اضافه می‌شود. مسیر‌هایی که گردیان از آن‌ها برمی‌گردد با نقطه‌چین قرمز مشخص شده‌اند
  			\cite{Wang2018LowShotLF}.
  		}	
  		\label{fig:imaginary1}
  	\end{center}
  \end{figure}  

  مقاله‌ی دیگری با الهام از انسان می‌خواهد وقتی یک وضعیت از یک شی را می‌بیند، بتواند آن تصویر را از زاویه‌های دیگر تصور کند و با دیدن یک نمونه، به توانایی شناخت برسد. در واقع به دنبال یادگیری مفاهیم است تا بتواند تبدیل‌های صحیح را پیاده کند. البته بیان می‌کند هر تبدیل معناداری ممکن است به بهتر شدن دقت دسته‌بندی شبکه‌ی عمیق کمک نکند. به همین دلیل، به دنبال تبدیل‌هایی است که تغییراتی به وجود آورند تا دقت مدل بیشتر شود
  \cite{Wang2018LowShotLF}
  . 

با روندی که در شکل 
\ref{fig:imaginary1}
مشخص شده، پیشنهاد داده‌اند که یک قسمت 
\trans{خیال‌باف}{Hallucinator}
با الگوریتم متایادگیری ترکیب شود و هر دو به صورت توأمان آموزش ببینند. به این منظور، از یک شبکه‌ی G
استفاده می‌کند که با ورودی گرفتن یک نمونه از داده‌های پشتیبان و یک نویز تصادفی، به صورت 
\trans{قطعی}{Deterministic}
یک داده‌ی جدید تولید کند. این داده‌های تولیدی را با نمونه‌های پشتیبان یک وظیفه ترکیب می‌کند و به وسیله‌ی هر دوی این‌ها یک طبقه‌بند h 
را آموزش می‌دهد. گرادیان تابع هزینه‌ی این دسته‌بند به G هم می‌رسد و به این ترتیب قسمت خیال‌باف یاد می‌گیرد تابع قطعی‌ای را تخمین بزند که به بهبود دقت دسته‌بند کمک کند. 

در کار یانگ و همکارانش، مشاهده شده آماره‌های دسته‌هایی که از لحاظ بصری به هم شبیه‌اند، مشابهت زیادی دارند
\cite{Yang2021FreeLF}
.
این روش روی مدل‌های مبتنی بر متر پیشنهاد شده است که برای هر دسته در فضا یک نماینده در نظر می‌گیرند و روی آن توزیعی متصور هستند. آن‌ها مشاهده کرده‌اند که آماره‌هایی مثل میانگین و پراکندگی توزیع دو دسته مثل روباه و گرگ که از لحاظ بصری به هم شبیه‌اند، حدود ۹۷ درصد مشابه است. در حالی که این حد از مشابهت با آماره‌های دسته‌های دیگر دیده نشد. 

آن‌ها فرض کرده اند که فضای بردار ویژگی‌ها حول هر دسته از توزیع 
\trans{گاوسی}{Gussian}
می‌آید. در نتیچه، به تعداد محدودی دسته‌ی با تعداد نمونه‌ی بالا را جمع‌اوری می‌کنند و آماره‌های مربوط به آن را حساب می‌کنند. سپس برای هر دسته‌ی نمونه محدود، نزدیک‌ترین دسته‌ها با تعداد نمونه‌‌ی بالا را پیدا می‌کنند. به دلیل مشابهتی که آماره‌های این دو دسته دارد، برای کلاس نمونه محدود آماره‌اش را 
\trans{تنظیم}{Calibrate}
می‌کنند و از آن توزیع نمونه‌برداری می‌کنند تا داده‌های بیشتری به مجموعه‌ی آموزش اضافه شود. روند این کار در شکل 
\ref{fig:calibration}
مشخص است.

\begin{figure}[!htbp]
	\begin{center}
		\includegraphics[width=1\textwidth]{images/calibration.png}
		%		\vspace{0.5in}
		\caption{
			قسمت خیال‌باف با G
			و دسته‌بند با h نشان داده شده است. برای یک اپیزود متایادگیری، به نمونه‌های پشتیبان هر دسته تعدادی داده اضافه می‌شود. مسیر‌هایی که گردیان از آن‌ها برمی‌گردد با نقطه‌چین قرمز مشخص شده‌اند
			\cite{Yang2021FreeLF}.
		}	
		\label{fig:calibration}
	\end{center}
\end{figure}  
  

\subsection{داده‌افزایی در سطح نمونه‌}

منظور از داده‌افزایی در سطح نمونه اعمال تبدیل‌هایی در فضای تصویر‌های ورودی به شبکه‌ی عمیق است
\cite{Shorten2019ASO}
. 
به این منظور از تبدیل‌هایی نظیر تغییرات در فضای رنگ، پاک کردن تصادفی بخشی از تصویر، ترکیب بخش‌های مختلف از دو تصویر، چرخش تصویر، 
\trans{انتقال سبک با شبکه‌‌های ژرف}{Neural Style Transfer}
و تولید داده با مدل‌های مولد برای داده‌افزایی استفاده می‌شود
\cite{Gatys2015ANA,Zhong2017RandomED,Takahashi2018DataAU}
.

\section{متایادگیری}

اولین واژه‌ای که برای متایادگیری میتوان متصور شد، یادگیری برای یادگیری است. اگر روی یادگیری انسانها تمرکز کنیم، نوعی استفاده از تجربیات پیشین برای یادگیری یک وظیفه‌ی جدید در رفتار آن‌ها دیده می‌شود. این مسئله یادگیری را برای انسان بسیار داده کارآمد می‌کند. در واقع انسان‌ها به عنوان موجوداتی هوشمند، از تجربیات پیشین در یادگیری وظیفه‌ها به نوعی استفاده می‌کنند که با کشف 
\trans{ساختارهای مشترک}{Shared Structures}
، می‌توانند عملکرد بهینه در یادگیری وظیفه‌ی جدید داشته باشند. 

انسان برای یکبار قوانین فیزیک حاکم بر محیط اطرافش را یاد میگیرد و هر بار در انجام یک فعالیت فیزیکی الگویی مشابه را می‌یابد، با الهام گرفتن از تجربیات پیشین، آن مسئله را حل میکند. وقتی یک فرد ورزشی مثل والیبال را تمرین می‌کند، در ورزشی مانند بسکتبال، در مقایسه با فردی که تجربه‌ی والیبال بازی کردن را ندارد، معمولا سریع‌تر یاد می‌گیرد. این به دلیل یادگرفتن چگونگی یادگیری در شرایطی با ساختار مشابه توسط انسان است.
در شبکه‌های عصبی، معنای ساده‌ی یادگیری این است که پارامترهای
\trans{بهینه}{Optimum}
 برای حل یک وظیفه مشخص شود. اما وقتی میخواهیم یادگیری برای یادگیری انجام دهیم، یعنی شبکه توانایی چگونگی بهینه یاد گرفتن را پیدا کند. در این صورت ویژگی‌هایی مانند داده کارآمد بودن،
 \trans{تطبیق سریع}{Early Adaptation}
  و تعمیم‌پذیری بالا به شبکه اضافه می‌شود.
  
  ما در مسئله‌ی متایادگیری، با دسته‌ای از مسائل یادگیری مواجه هستیم که هر کدام از آن‌ها را یک وظیفه می‌نامیم و می‌خواهیم در طول یادگرفتن وظایف مختلف، خود فرآیند یادگیری را بهتر کنیم. در زیر مجموعه وظایف یادگیری در قالب 
  $\mathcal{D}_{meta}$
  نشان داده شده است.
  \begin{equation}
  	\begin{split}
  		\mathcal{D}_{meta} = \{\mathcal{T}^{(1)}, \mathcal{T}^{(2)} , ..., \mathcal{T}^{(n)}\}
  		\end{split}
  \end{equation}

هر وظیفه‌ی یادگیری برای هر دسته، داده‌های آموزش دارد که به آن‌ها داده‌‌های پشتیبان می‌گوییم و با 
$\mathcal{S}$
نشان می‌دهیم. همچنین به داده‌های آزمون، داده‌ی پرسش می‌گوییم و با 
$\mathcal{Q}$
نمایش می‌دهیم.
\begin{equation}
	\begin{split}
		\mathcal{T}^{(i)} = \{S^{(i)},Q^{(i)}\}
	\end{split}
\end{equation}

در مسائل متایادگیری، دو دسته پارامتر وجود دارد. پارامتر‌های 
$\theta$
که به آن‌ها متاپارامتر می‌گوییم و پارامتر‌‌های کلی مدل هستند و برای تمام وظایف و مسائل در 
  $\mathcal{D}_{meta}$
  مشترک‌اند. دسته‌ی دیگر پارامتر‌های 
  \trans{خاص وظیفه}{Task Specific}
  هستند و همانطور که از اسم آن پیداست، مختص هر یک از وظایف می‌باشند. این دسته از پارامتر‌ها را معمولا با 
  $\phi$
  در ادبیات متایادگیری نمایش می‌دهند. 
  
  در واقع متاپارامتر‌ها همان پارامتر‌هایی هستند که تجربیات مشترک حاصل از دیدن مسائل و وظایف مختلف یادگیری را در خود جای داده‌اند. با یادگرفتن متاپارامترهاست که فرآیندهای یادگیری جدید بهینه می‌شود و وقتی متاپارامترها آموزش داده می‌شوند، مدل یاد می‌گیرد که چگونه یاد بگیرد. اما 
  $\phi$
در قالب یکی از وظیفه‌ها یا همان یکی از مسائل یادگیری تعریف می‌شود. 

در واقع پارامتر‌های خاص وظیفه، برای یک وظیفه‌ی یادگیری 
$i$
، از روی 
$S^{(i)}$
و با کمک و استفاده از اطلاعات متاپارامترها، یعنی 
$\theta$
یاد گرفته می‌شوند. بعد از اینکه 
  $\phi$
  یاد گرفته شد، به سراغ داده‌های آزمون یا پرسش یک وظیفه یعنی 
  $Q^{(i)}$
  می‌رویم و با استفاده از آن‌ها، از مدل استنتاج می‌گیریم و به وسیله‌ی تعریف یک تابع هزینه، خطایی که بدست می‌آید را برای به روزرسانی متاپارامترها استفاده می‌کنیم. با تکرار این روند در طول وظیفه‌‌های مختلف، متاپارامترهای
  $\theta$
  به مقدار مناسبی همگرا می‌شوند به نحوی که از روی آن‌ها و داده‌های پشتیبان، با سرعت بالا و با دیدن نمونه‌‌های کم می‌توان پارامتر‌های خاص وظیفه را بدست آورد. به فرآیند یادگیری متاپارامترها که در طول وظایف مختلف صورت می‌گیرد، 
  \trans{حلقه‌ی بیرونی}{Outer Loop}
   یادگیری و به مرحله‌ی بهینه‌سازی برای بدست آمدن پارامتر‌های خاص یک وظیفه، 
   \trans{حلقه‌ی درونی}{Inner Loop}
   یادگیری یا مرحله‌ی 
   \trans{یادگیری سریع}{Fast Learning}
    می‌گویند. در شکل 
    \ref{fig:meta1}
    کلیت این روش آمده است.
\begin{figure}[!htbp]
	\begin{center}
		\includegraphics[width=1\textwidth]{images/meta1.png}
		%		\vspace{0.5in}
		\caption{
			آموزش به شیوه‌ی متایادگیری
			\cite{Ravi2016OptimizationAA}.
		}	
		\label{fig:meta1}
	\end{center}
\end{figure}  

همانطور که در یک مسئله‌ی یادگیری، برای سنجیدن کیفیت پارامترهای یادگرفته شده، از مرحله‌ی آزمون و مجموعه‌ی دادگان آزمون استفاده می‌کنیم، برای سنجیدن کیفیت متاپارامتر‌های یادگرفته شده هم از یک فاز متاآزمون مطابق شکل 
\ref{fig:meta1}
استفاده می‌کنیم. 

برای این منظور، وظایف جدیدی می‌سازیم که دسته‌های موجود در آن‌ها، در مرحله‌‌ی قبلی یعنی متایادگیری وجود نداشته باشند. یعنی دسته‌های موجود در مرحله‌ی متاآزمون با دسته‌های موجود در مرحله‌ی متایادگیری 
\trans{ناسازگار}{Disjoint}
اند
(برای یک وظیفه در مرحله‌ی متاآزمون 
$\mathcal{T} \notin \mathcal{D}_{meta}$)
.
به طور خلاصه، الگوریتم متایادگیری از دو مرحله تشکیل شده است. در مرحله‌ی متایادگیری مدل به دنبال بهینه‌سازی و یادگیری متاپارامترها 
$\theta$
است و می‌خواهد تجربیات حاصل از مسائل این مرحله را در خود جای دهد. این دانش پیشین در مرحله‌ی متاآزمون مورد بهره‌برداری قرار می‌گیرد و به کمک آن مدل باید پارامتر‌های خاص هر وظیفه‌ی جدید را با دیدن نمونه‌های محدود به خوبی بدست آورد. در ادامه، ابتدا نمادگذاری دقیق‌تر مسئله‌ی متایادگیری را بیان می‌کنیم و سپس به معرفی تقسیم‌بندی مسائل در این حیطه می‌پردازیم
\cite{Hospedales2020MetaLearningIN}
. 

\subsection{ساختار کلی الگوریتم متایادگیری}

معمولا الگوریتم یادگیری که برای یادگیری پارامتر‌های خاص تسک استفاده می‌شود را با 
$\mathcal{B}$
نشان می‌دهند. در مرحله‌ی یادگیری سریع یا حلقه‌ی داخلی، یک مدل یا به طور معادل الگوریتم یادگیری را انتخاب می‌کنیم که به آن 
\trans{یادگیر پایه}{Base Learner}
می‌گوییم. برای آموزش این یادگیر پایه، با استفاده از داده‌های پشتیبان یک وظیفه و همینطور دسترسی به متاپارامترها یک مسئله‌ی بهینه‌سازی به صورت زیر را حل می‌کنیم.
\begin{eqnarray}
	\phi = \mathcal{B}\left(\mathcal{S} ; \theta \right)=\underset{\phi}{\arg \min } \ \mathcal{L}^{{base}}\left(\mathcal{S} ; \theta, \phi\right)
\end{eqnarray} 

در این معادله، منظور از 
$\mathcal{L}^{base}$
تابع هزینه‌ی مختص یک وظیفه است. با بدست آمدن پارامترهای 
  $\phi$
  که در واقع وزن‌های یادگیر پایه اند، یک مدل
  $\mathcal{M}$
  که نشان‌دهنده‌ی متایادگیر است آموزش می‌دهیم. برای این منظور مسئله‌ی بهینه‌سازی زیر را حل می‌کنیم.
    \begin{equation}
    	\theta^{\star} = \mathcal{M}(\mathcal{D}_{meta})= \underset{\theta}{\arg \min} \  \mathbb{E}_{(\mathcal{S}, \mathcal{Q})\in \mathcal{D}_{meta}}[\mathcal{L}^{{meta }}\left(\mathcal{Q}; \theta, \phi\right), \ where \ \phi=\mathcal{B}(\mathcal{S}; \theta)]
    \end{equation}
در معادله‌ی بالا منظور از 
$\mathcal{L}^{meta}$
تابع هزینه‌ی است که مختص ارزیابی متاپارامترها می‌باشد. در ادبیات متایادگیری، به در محدوده‌ی یک وظیفه، به تعداد نمونه‌های پشتیبان به ازای هر دسته تعداد 
\trans{شات}{Shot}
و به تعداد دسته‌های موجود 
\trans{راه}{Way}
می‌گویند.

%\subsection{انواع روش‌های متداول متایادگیری}

در رویکرد متایادگیری، انواع پیاده‌سازی‌های متفاوت از این مفهوم وجود دارد. این پیاده‌سازی‌ها گستره‌ای از روش‌های مختلف را در بر می‌گیرد که هر کدام ویژگی‌های خاصی دارند و در ادامه به معرفی آن‌ها می‌پردازیم.

\subsection{روش‌های مبتنی بر بهینه‌سازی}

در این دسته از روش‌ها در یک فرآیند بهینه‌سازی دو سطحی، در فاز متاآموزش، در هر وظیفه پارامترهای خاص آن وظیفه را یاد می‌گیریم. سپس با بدست آوردن خطای داده‌های پرسش، بهینه‌سازی روی متاپارامترها صورت میگیرد.

مقاله‌ی MAML
به عنوان شاخص‌ترین روش این دسته مطرح است
\cite{MAML}
.
تابع هدف حاصل از این روش در معادله‌ی
\ref{eq:maml}
 آمده است. در این معادله،
 $\mathcal{T}_i$
  یک وظیفه‌ی نمونه‌برداری شده از توزیع مجموعه وظایف مسئله‌ی هدف است. تابع خطا با
  $\mathcal{L}$
   مشخص شده و از
  $\alpha$
    برای نمایش نرخ یادگیری استفاده شده است. همچنین 
    $\mathcal{D}_i^{support}$
    مجموعه‌ی داده‌‌های پشتیبان برای یک وظیفه و 
    $\mathcal{D}_i^{query}$
    مجموعه‌ی داده‌های پرسش برای وظیفه‌ی i است. همانطور که قبلا بیان شد، 
    $\theta$
    متاپارامترها را نشان می‌دهد و منظور از 
    $\theta^{'}$
    پارامتر‌های خاص وظیفه یا همان 
    $\phi$
    است که در این روش با بهینه‌سازی و آپدیت متاپارمتر‌ها بدست می‌آید.
    \begin{equation}
    	\label{eq:maml}
    	\begin{split}
    		\min_{\theta} \sum_{\mathcal{T}_{i} \sim p(\mathcal{T})} \mathcal{L}_{\mathcal{T}_{i}}(f_{\theta^{'}_{i}}) = \sum_{\mathcal{T}_{i} \sim p(\mathcal{T})} \mathcal{L}(\theta - \alpha \nabla_{\theta}\mathcal{L}(\theta, \mathcal{D}_i^{support}),\mathcal{D}_i^{query})
    \end{split}\end{equation}

در شکل 
\ref{fig:maml}
شهود این روش آمده است.

\begin{figure}[!htbp]
	\begin{center}
		\includegraphics[width=1\textwidth]{images/maml.png}
		%		\vspace{0.5in}
		\caption{
			متایادگیری مستقل از مدل؛ در این روش، با تعداد کمی بهینه‌سازی روی متاپارامترها، به پارامترهای خاص هر وظیفه می‌رسیم
			\cite{MAML}.
		}	
		\label{fig:maml}
	\end{center}
\end{figure}  

در واقع متاپارامتر‌ها یک نقطه‌ی شروع اولیه برای پارامترهای خاص هر وظیفه اند و با داشتن اطلاعات مشترک بین وظایف، سبب می‌شوند با تعداد کمی گام بهینه‌سازی، به پارامتر‌های مناسب برای حل یک وظیفه برسیم.  این روش قادر به یادگیری توابع با پیچیدگی بالا است اما به دلیل نیاز به
\trans{بهینه‌سازی دو سطحی}{Second Order Optimization}
 و ظهور مشتق دوم، مشکلاتی از لحاظ حافظه به جهت نیاز به نگه داشتن تمامی مشتق‌های مرتبه  اول و همینطور پیچیدگی فرآیند یادگیری را به دنبال دارد. برای حل این مشکلات راهکارهایی از قبیل استفاده از تقریب مرتبه‌اول مشتق‌گیری و همچنین بهینه‌سازی محدب با فرم بسته پیشنهاد شده است
\cite{Nichol2018OnFM,Lee2019MetaLearningWD}
.

\subsection{روش‌های 
\trans{جعبه سیاه}{Black Box}
}

در این دسته از روش‌ها، یک مدل جداگانه آموزش می‌دهیم که با دریافت توصیف هر وظیفه، بخشی یا تمامی پارامتر‌های مدل مربوط به آن وظیفه را خروجی بدهد. در واقع یک مدل داریم با پارامترهای 
$\theta$
که به عنوان ورودی توصیف یک وظیفه یا همان داده‌های پشتیبان را می‌گیرد و سپس در خروجی یک 
\trans{آماره‌ی کافی}{Sufficient Statistic}
از پارامتر‌های خاص تسک یا همان 
$\phi$
را تولید می‌کند.
\begin{figure}[!htbp]
	\begin{center}
		\includegraphics[width=1\textwidth]{images/snail.png}
		%		\vspace{0.5in}
		\caption{
			متایادگیری جعبه سیاه؛ در شکل سمت چپ که تنظیمات یادگیری با ناظر آمده است، می‌توان از 
			$X_{t-3}$
			تا 
			$X_{t-1}$
			را به همراه برچسب هر کدام به عنوان مجموعه‌ی داده‌های پشتیبان در نظر گرفت و 
			$X_t$
			را داده‌ی پرسش در نظر گرفت. به دلیل ماهیت شبکه‌های بازگشتی، با به روز رسانی فضای نهان، به پارامتر‌های خاص وظیفه می‌رسیم
			\cite{Mishra2017ASN}.
		}	
		\label{fig:snail}
	\end{center}
\end{figure}  

بخش‌های اساسی این دسته از مدل‌ها، 
$g_{\phi}$
یا همان مدل خاص وظیفه و 
$f_{\theta}$
یا همان مدل جعبه سیاه است که مدل خاص وظیفه را با دریافت داده‌های پشتیبان یک وظیفه می‌سازد. عملکرد این دو مدل مطابق فرمول 
\ref{eq:snail}
است.
\begin{eqnarray}
	\label{eq:snail}
	\begin{array}{c}
		\phi = \mathcal{B}\left(\mathcal{S} ; \theta \right)= f_\theta(\mathcal{S}) \\
		\hat{y} = g_\phi(x)		
	\end{array}
\end{eqnarray}

معماری‌های رایج در این دسته شامل 
\trans{مجموعه‌ی عمیق}{Deep Set}
\cite{DeepSet}
، 
\trans{مکانیزم توجه}{Attention Mechanism}
\cite{Attention}
و 
\trans{شبکه‌‌های بازگشتی}{Recurrent Neural Networks}
\cite{LSTM}
است. این دسته از مدل‌ها در صورت بزرگ بودن تعداد پارامتر‌های مدل خاص وظیفه، عملکرد مناسبی ندارند.

\subsection{روش‌های 
\trans{مبتنی بر متر}{Metric Based}
}

با توجه به پیچیدگی‌های روش دسته‌های قبل در تخمین پارامترهای خاص وظیفه، روش مبتنی بر متر یا بدون پارامتر پیشنهاد شده است. در این دسته از روش‌ها، سعی می‌شود در یک فضای معنایی، مبتنی بر یک معیار فاصله دسته‌بندی انجام پذیرد. نکته‌ی حائز اهمیت این است که مرحله‌ی متاآموزش همچنان دارای پارامتر است و در واقع در این مرحله متاپارامترها یاد گرفته می‌شوند. اما در فاز استنتاج، دسته‌بندی مبتنی بر متر و فاصله انجام میگیرد. تمامی کارهای پیشین این دسته بر پایه‌ی شبکه‌ی مبتنی بر 
\trans{نماینده}{Prototype}
 است
 \cite{protonet}
 .
 \begin{figure}[!htbp]
 	\begin{center}
 		\includegraphics[width=1\textwidth]{images/protonet.png}
 		%		\vspace{0.5in}
 		\caption{
 			متایادگیری مبتنی بر متر؛ در این دسته از روش‌ها برای هر دسته یک مرکز از روی نمونه‌های پشتیبان هر دسته یا متادیتای آن محاسبه می‌شود و برای یک نمونه پرسش، با سنجش فاصله از مرکز دسته‌ها، طبقه‌بندی انجام می‌شود. در سمت چپ محاسبه‌ی نماینده دسته از روی نمونه‌‌های پشتیبان و در سمت راست محاسبه‌ی نماینده دسته از روی متادیتای دسته‌ها انجام شده است
 			\cite{protonet}.
 		}	
 		\label{fig:protonet}
 	\end{center}
 \end{figure}  
 
 در این روش، با همان تنظیماتی که به طور کلی برای متایادگیری تعریف کردیم، ابتدا یک شبکه‌ی نگاشت فضای ویژگی‌ها آموزش داده می‌شود. در واقع پارامترهای این شبکه در طول وظیفه‌های مختلف یاد می‌گیرند چگونه یک داده‌ی ورودی را به فضایی نگاشت کنند که در آن دوری و نزدیکی نقاط از هم معنا داشته باشد. پارامتر‌های این شبکه 
 $\theta$
 را تشکیل می‌دهند. در چنین فضایی، وقتی دو داده از یک دسته نگاشت می‌شوند، باید فاصله‌ی آنها نزدیک به هم باشد. به طور قرینه، اگر دو داده از دسته‌های متفاوتی باشند، باید به نقاطی دور از هم نگاشت شوند. با پیاده‌سازی این ایده، اگر برای هر دسته یک نماینده در فضای ویژگی‌ها محاسبه کنیم، در صورت دیدن یک داده‌ی مجهول، می‌توان با استفاده از شبکه‌ی آموزش دیده، نگاشت روی فضای معنایی را انجام داد و سپس با یک معیار فاصله مانند فاصله‌ی اقلیدسی، تعلق داده به نزدیکترین نماینده را مشخص کرد. در واقع متاپارامترها در این مسئله، همان پارامترهای شبکه‌ی نگاشت به فضای معنایی است.
 
 \begin{equation}
 	\label{eq:proto}
 	\begin{split}
 		c_{k}=\frac{1}{|S_{k}|}\sum_{(x_{i},y_{i}) \in{k}} f_{\theta}(x) \\
 		 p_{\theta}(y=k|x)=\frac{\exp(-d(f_{\theta}(x), c_{k}))}{\sum_{k^{'}} \exp(-d(f_{\theta}(x), c_{k^{'}})} \\ min J(\theta)=-\log p_{\theta}(y=k|x)
 	\end{split}
 \end{equation}

	در معادله‌ی 
\ref{eq:proto}
و برای دسته‌ی 
$k$
که به تعداد 
$|S_k|$
عضو دارد، 
ابتدا نماینده هر دسته توسط شبکه‌ی نگاشت به فضای ویژگی 
$f$
محاسبه و در 
$c_k$
ذخیره شده است. سپس با یک معیار فاصله‌ی 
$d$
نزدیکی یک داده‌ی ورودی با هر یک از نماینده دسته‌ها محاسبه شده است و احتمال تعلق داده به هر دسته بدست آمده است.
تابع هدف هدف این روش پیدا کردن متاپارامترهای 
$\theta$
به منظور کمینه‌کردن خطای
\trans{منفی بیشینه‌‌ی درست‌نمایی}{negative log-likelihood} 
تعریف می‌شود که با 
$J$
نشان داده شده است.

\begin{latin}
	\begin{algorithm}
		\caption{Training episode loss computation for prototypical networks. N is the number of
			examples in the training set, K is the number of classes in the training set, $N_C \le K$ is the number
			of classes per episode, $N_S$ is the number of support examples per class, $N_Q$ is the number of query
			examples per class. \text{RandomSample}$(S, N)$ denotes a set of N elements chosen uniformly at
			random from set S, without replacement.}
		\label{alg:protonet}
		\begin{algorithmic}[1]
			\renewcommand{\algorithmicrequire}{\textbf{Input:}}
			\renewcommand{\algorithmicensure}{\textbf{Output:}}
			\renewcommand{\algorithmiccomment}[1]{\hfill \textit{//#1}}
			\REQUIRE Training set $D = \{(x_1, y_1), \ldots, (x_N, y_N)\}$, where each $y_i \in \{1, \ldots, K\}$. $D_k$ denotes the subset of $D$ containing all elements $(x_i, y_i)$ such that $y_i = k$.
			
			\ENSURE The loss $J$ for a randomly generated training episode.
			
			\STATE $V \leftarrow \text{RandomSample}(\{1, \ldots, K\}, N_C)$
			\hfill \textit{//Select class indices for episode}
			
			\FOR{$k$ in $\{1, \ldots, N_C\}$}
			\STATE $S_k \leftarrow \text{RandomSample}(D_{V_k}, N_S)$
			\hfill \textit{// Select support examples}
			\STATE $Q_k \leftarrow \text{RandomSample}(D_{V_k} \setminus S_k, N_Q)$
			\hfill \textit{//Select query examples}
			\STATE $c_k \leftarrow \frac{1}{N_C} \sum_{(x_i, y_i) \in S_k} f_{\phi}(x_i)$
			\hfill \textit{//Compute prototype from support examples}
			\ENDFOR
			
			\STATE $J \leftarrow 0$
			\hfill \textit{//Initialize loss}
			
			\FOR{$k$ in $\{1, \ldots, N_C\}$}
			\FOR{$(x, y)$ in $Q_k$}
			\STATE $J \leftarrow J + \frac{1}{N_C \cdot N_Q} \left( d(f_{\phi}(x), c_k) + \log \sum_{k'} \exp(-d(f_{\phi}(x), c_{k'})) \right)$
			\hfill \textit{//Update loss}
			\ENDFOR
			\ENDFOR
		\end{algorithmic}
	\end{algorithm}
\end{latin}


انتخاب تابع فاصله‌ی مناسب یکی از چالش‌های این روش است. در اغلب مواقع تابع فاصله، یک تابع بدون پارامتر و ساده است
\cite{Liu2020FewShotOR}
.
اما در برخی پژوهش‌ها، از تابع‌های فاصله‌ی پارامتری ‌بسته به کاربرد مسئله، بهره برده اند
\cite{Sung2017LearningTC}
.
مزیت این روش‌ها علاوه بر سادگی و سرعت بالا در فاز استنتاج، نتایج بهتر آن نسبت به روش‌های دیگر متایادگیری است. به طور عمده در آزمایش‌های مختلف نشان داده شده است روش‌های مبتنی بر نماینده بر روش‌های دیگر برتری دارد. به این جهت، می‌توان دریافت پیدا کردن یک نگاشت مناسب به فضای معنایی ویژگی‌ها که در آن یک فاصله مناسب جهت دسته‌بندی وجود دارد، برای حل مسائل بزرگ کافی است.

از مزایای این روش، بدون پارامتری بودن مرحله‌ی استنتاج است. این موضوع ما را از بدست آوردن پارامتر‌های خاص هر وظیفه به صورت یک فرآیند یادگیری بی‌نیاز می‌کند و دسته‌بندی با رویکرد مولد انجام می‌شود. به همین دلیل اگر دسته‌ی جدیدی در فاز آزمون اضافه شود، نیاز به به روز رسانی پارامتر‌ها نیست و فقط کافی است نماینده‌ی آن دسته محاسبه شود. همچنین این دسته از روش‌های متایادگیری در مقابل سایر رویکرد‌ها به دقت‌های بالاتری در مسائل گوناگون رسیده است. در پژوهش حاضر، از این رویکرد متایادگیری استفاده کرده‌ایم. به همین جهت نیاز به دقت بیشتر در شیوه‌ی یادگیری آن داریم که به صورت شبه کد در الگوریتم
\ref{alg:protonet}
 آمده است.

%\begin{latin}
%	\begin{algorithm}
%		\caption{Coarse-grained classifier (training)}\label{alg:coarse:train}
%		\begin{algorithmic}[1]
%			 \renewcommand{\algorithmicrequire}{\textbf{Input:}}
%			\renewcommand{\algorithmicensure}{\textbf{Output:}}
%			\renewcommand{\algorithmiccomment}[1]{\hfill \textit{//#1}}
%			\REQUIRE dataset $\mathcal{D}_{source}$ and $\mathcal{D}_{target}$, backbone $g_\phi$, classification head $f_\varphi$
%			
%			\ENSURE backbone $g_{\phi'}$, classification head $f_{\varphi'}$
%			
%			\WHILE[Supervised learning]{not converged}
%			\STATE sample batch $\{(x^i, y_{super}^i)\}_{i=1}^B\sim\mathcal{D}_{source}$ 
%			\STATE $\hat{y}_{super}^i \leftarrow f_\varphi(g_\phi(x^i))$
%			\hfill \textit{//Do this for all of the batch samples}
%			\STATE update $\phi$ and $\varphi$ based on $\frac{1}{B}\sum\limits_{i=1}^B\mathcal{L}(y^i_{super}, \hat{y}^i_{super})$
%			\ENDWHILE
%			
%			\STATE initialize $\phi' \leftarrow \phi, \varphi' \leftarrow \varphi$
%			
%			\FOR[Semi-supervised learning]{$c \in \mathcal{Y}_{base, target}$}
%			\STATE take all support samples $\mathcal{S}^c$
%			\STATE $\Tilde{y}_{super} \leftarrow \operatorname*{argmax} \prod\limits_{x^i \in \mathcal{S}^c} f_{\varphi'}(g_{\phi'}(x^i))$
%			\STATE $\hat{y}_{super}^i \leftarrow f_{\varphi'}\left(g_{\phi'}\left(x^i\right)\right)$
%			\hfill\textit{//Do this for all of the support samples}
%			\STATE update $\phi'$ and $\varphi'$ based on $\frac{1}{|\mathcal{S}^c|}\sum\limits_{i=1}^{|\mathcal{S}^c|}\mathcal{L}(\Tilde{y}_{super}, \hat{y}^i_{super})$
%			\ENDFOR
%		\end{algorithmic}
%	\end{algorithm}
%\end{latin}


\subsection{اشتراک تجربه‌ی احتمالاتی}

در بررسی شبکه‌های متایادگیری، یک نقطه‌ی ضعف مشترک بین روش‌های متایادگیری دیده شد. تمامی این شبکه‌ها فرض می‌کردند اطلاعات موجود در متاپارامترها باید بین همه‌ی وظیفه‌ها به اشتراک گذاشته شود
\cite{Vinyals2016MatchingNF,protonet,MAML}
.
این در حالی است که الگوی یادگیری انسان، تجربیات حاصل از تجربه روی وظیفه‌های مختلف را با همه‌ی وظیفه‌ها به اشتراک نمی‌گذارد. در واقع عملکرد به نحوی است که ابتدا یک مشابهت‌یابی بین وظیفه‌ی در حال انجام با تجربیات قبلی صورت می‌پذیرد و از هر کدام از آن‌ها بخشی از تجربیات مرتبط به کمک یادگیری بهتر و سریع‌تر وظیفه‌ی فعلی می‌آید. جهت اصلاح این روند، کار‌هایی جهت اصلاح روند دانش قابل انتقال به هر وظیفه صورت پذیرفته است. یک روند رویکرد متایادگیری احتمالاتی بوده است که برای هر وظیفه، یک مدل از توزیع تمامی مدل‌های ممکن نمونه گرفته می‌شود
\cite{Lee2018GradientBasedMW, Kim2018BayesianMM, Finn2018ProbabilisticMM}
.
در این طیف روش‌ها، تمرکز بر روی انتقال دانش به وظیفه هدف است. به همین دلیل تعمیم‌پذیری دانش بین وظیفه‌های با 
\trans{همبستگی بالا}{Highly Correlated}
اتفاق نمی‌افتد. در واقع در لحظه‌ی اشتراک دانش، فقط به وظیفه‌ی فعلی توجه می‌شود و ارتباط آن با وظایف دیگر در نظر گرفته نمی‌شود. 
\begin{figure}[!htbp]
	\centering
	\includegraphics[width=0.85\textwidth]{images/sharemeta.png}
	\caption{انواع مختلف انتقال تجربه از متاپارامتر‌ها به پارامتر‌های خاص وظیفه؛ در شکل سمت چپ، متاپارامتر‌ها به صورت یکسان در اختیار تمامی وظیفه‌ها قرار می‌گیرد. در شکل وسط، به ازای هر وظیفه خاص، یک متاپارامتر خاص نمونه‌گیری می‌شود و به آن وظیفه انتقال می‌یابد. در شکل سمت راست، وظایف خوشه‌بندی می‌شوند و به هر خوشه، تجربیات یکسانی انتقال می‌یابد.
		\cite{Yao2019HierarchicallySM}}	
	\label{fig:sharemeta}
\end{figure}

برای بهبود این مسئله، باید به یک ساختار متایادگیری رسید که بتوان تعمیم‌پذیری و شخصی‌سازی دانش را برای هر وظیفه به انجام رساند. به این جهت در دسته‌ای دیگر از روش‌های پیشنهادی، با الهام از مفاهیم روانشناسی، هنگام مواجهه با یک وظیفه انسان به خوشه‌بندی
آن می‌پردازد و به این وسیله، شباهت وظیفه با تجربیات گذشته سنجیده می‌شود. در نتیجه انتقال تجربه برای یک وظیفه، از خوشه‌ی مرتبط با آن نشئت می‌گیرد. البته برخی از تجربیات هم به صورت 
\trans{جهانی}{Globallly}
به همه‌ی وظایف انتقال می‌یابند
\cite{Yao2019HierarchicallySM}
.
در شکل 
\ref{fig:sharemeta}
به ترتیب از چپ به راست سه رویکرد مطرح شده جهت استفاده از متاپارامتر‌ها برای رسیدن به پارامتر‌های خاص وظیفه مشخص شده است.

\section{استفاده از 
\trans{بدنه}{Backbone}
\trans{بیانگر}{Expressive}
 برای رسیدن به 
 \trans{کارایی نمونه}{Sample Efficiency}
}

در تنظیمات نمونه محدود، داده کارآمدی مدل اهمیت به‌سزایی دارد. یعنی مدل برای یادگرفتن از حداقل تعداد نمونه به ازای هر دسته استفاده کند. به این منظور، یک رویکرد مهم در پژوهش‌های اخیر و گذشته استفاده از 
\trans{یادگیری انتقالی}{Transfer Learninig}
و استفاده از 
\trans{بدنه}{Backbone}
های بزرگ است که بازنمایی‌های بیانگر خروجی می‌دهند. این بدنه‌ها معمولا بر روی مجموعه دادگان بزرگ از قبل پیش‌آموزش دیده‌اند
\cite{Hu2022PushingTL,Tian2020RethinkingFI,Xu2021UnsupervisedMF}
.

علاوه بر یادگیری نظارت‌شده، یادگیری بدون ناظر
\cite{Tian2020RethinkingFI,Xu2021UnsupervisedMF}
و یادگیری خودنظارتی
\cite{Lu2022SelfSupervisionCB}
 نشان داده‌اند که می‌توانند روش‌های مناسبی برای آموزش یک بدنه‌ی مناسب برای یادگیری نمونه محدود باشند. اخیرا، مدل‌های بنیادی بزرگ مانند 
 CLIP
 \cite{clip}
 و 
 DINO
 \cite{dinov1}
 با استقبال گسترده‌ای در حیطه‌ی یادگیری نمونه محدود و 
 \trans{بدون نمونه}{Zero Shot}
  مواجه شده‌اند
  \cite{Zhang2023PromptGT,Zhang2022TipAdapterTA}
  .
  در شدید‌ترین حالت یادگیری، هیچ نمونه‌ای برای یادگیری در اختیار نداریم و باید برای یک داده‌ی آزمون، استنتاج را انجام دهیم. در این حالت نیاز به ترکیب اطلاعات زبانی و تصویری داریم. به نحوی که بازنمایی برچسب‌های هر دسته، در جایی از فضا نگاشت شوند که نزدیک به بازنمایی تصاویر همان دسته باشد. برای حل این دسته از مسائل، از یادگیری تباینی استفاده می‌شود و مهم‌ترین معماری آن CLIP 
  است که در شکل 
  \ref{fig:clip}
  آمده است.
  
 \begin{figure}[!htbp]
	\begin{center}
		\includegraphics[width=1\textwidth]{images/clip.png}
		%		\vspace{0.5in}
		\caption{
			مدل CLIP
			سعی می‌کند بازنمایی تصاویر و برچسب‌های متناظر که روی قطر اصلی قرار دارند را به هم نزدیک کند. در نهایت از این ویژگی برای دسته‌بندی بدون نمونه استفاده می‌کند
			\cite{clip}.
		}	
		\label{fig:clip}
	\end{center}
\end{figure}  

\section{خاصیت مجموعه باز بودن}

برای رویارویی با مجموعه باز بودن دسته‌بندی دو رویکرد اساسی پیشنهاد شده است. یک رویکرد سلبی، دادن
\trans{خودآگاهی}{Self Awareness}
 به مدل است. یعنی یک مدل آگاهی این را داشته باشد که روی چه ورودی می‌تواند خروجی تولید کند. در نتیجه وقتی یک ورودی خارج از
حیطه‌ی توانایی‌های مدل است، از اظهار نظر روی آن خودداری کند. البته داشتن خودآگاهی، گروه بزرگتری از مسائل را در بر می‌گیرد که به سه دسته‌ی عمده تقسیم می‌شوند. گروه اول تشخیص 
\trans{خارج از توزیع}{Out of Distribution}
است. در این گروه تصاویری که از توزیع دادگان آموزش نیستند یا از دامنه‌های دیگر مانند
\trans{حملات تخاصمی}{Adverserial Attacks}
 می‌آیند شناسایی می‌شود
 \cite{Hendrycks2016ABF,Huang2017AdversarialAO}
 .
 گروه دوم طبقه‌بندی
 \trans{واقع‌بینانه}{Realistic}
  است که وقتی یک نمونه برای طبقه‌بندی سخت به نظر می‌رسد، از دسته‌بندی آن خودداری می‌شود
  \cite{Lu2017SafetyNetDA}
  .
  گروه آخر هم مسئله‌ی شناسایی مجموعه باز است که وقتی نمونه‌ای از کلاس بدیع می‌بیند، آن را طبقه‌بندی نکند.
  

اما رویکرد دیگر که به نوعی ایجابی است، در حیطه‌ای است که
مدل، خود را با ذات مجموعه باز بودن مسئله تطبیق دهد و برای تعمیم دانسته‌هایش روی دسته‌های جدید آماده باشد. اگر مدل مولد باشد با تنظیم توزیعی که یادگرفته میتواند به سرعت با اضافه شدن دسته‌های جدید انطباق پیدا کرد. این خاصیت در مدل‌های بدون پارامتر هم نمود دارد. اما در مدل‌های تمایزگذار با اضافه شدن دسته‌های بدیع، نیاز به آموزش مدل از ابتدا بر روی تمامی دسته‌ها است. به همین جهت، در مسائل دسته‌بندی بزرگ مقیاس باید رویکردی را انتخاب کنیم که به کمترین میزان آموزش مجدد شبکه نیاز باشد. البته در مدل‌های تمایزگذار، روش‌هایی پیشنهاد شده است که بتوان وزن های مرتبط با دسته‌های جدید را تنها با طی کردن 
\trans{گذر جلو رو}{Forward Pass}
و بدون آموزش مجدد، بدست آورد.

\subsection{رویکرد سلبی؛ شناسایی دسته‌های بدیع}

در این رویکرد، سعی بر تشخیص دسته‌های بدیع و به نوعی عدم استنتاج مدل روی آن‌ها است. یعنی مدل فرآیندی برای تطابق با دسته‌های جدید و پذیرش آن‌ها طی نمی‌کند. عمده‌ی تمرکز پژوهش‌ها در این دسته روی تشخیص کلاس‌های بدیع در تنظیمات بزرگ مقیاس بوده است. اما در سال‌های اخیر روش‌هایی ارائه شده‌اند که علاوه بر مسائل بزرگ مقیاس با تعداد نمونه کافی، در مسائل نمونه محدود با تعداد دسته‌ی محدود هم عملکرد مناسبی داشته باشند.
در یک روش، به این نتیجه رسیده‌اند که دسته‌بندی مبتنی بر 
\trans{سافت‌مکس}{Softmax}
چون روی دسته‌های آموزش دچار 
\trans{بیش برازش}{Overfit}
است، نمی‌تواند راه حل خوبی برای تشخیص دسته بدیع باشد. به همین دلیل با الهام از رویکرد متایادگیری، در هر وظیفه دسته‌های دیده شده و بدیع را جای داده‌‌اند و علاوه بر تابع هزینه اصلی، یک جمله برای بیشینه کردن آنتروپی توزیع 
\trans{پسین}{Posterior}
برای دسته‌های بدیع قرار داده‌اند
\cite{Liu2020FewShotOR}
.

در یک کار دیگر، ابتدا یک 
\trans{خودکدگذار}{Auto Encoder}
به وسیله‌ی داده‌‌های درون توزیع آموزش داده می‌شود. بر روی این قسمت، یک لایه‌ی دسته‌بند هم اضافه می‌شود تا احتمال تعلق یک نمونه به دسته‌های درون توزیع را بدهد. سپس یک قسمت دیگر به وسیله‌ی داده‌های درون توزیع و دسته‌های بدیع آموزش داده می‌شود که باید تشخیص دهد آیا یک داده از کلاس بدیع هست یا خیر. عملکرد این مدل در شکل 
\ref{fig:openset1}
آمده است
\cite{Hendrycks2016ABF}
.
\begin{figure}[!htbp]
	\begin{center}
		\includegraphics[width=0.7\textwidth]{images/openset1.png}
%		\vspace{0.5in}
		\caption{
		استفاده از یک خودکدگذار و تولید احتمالات دسته‌بندی برای یک ورودی که بر روی داده‌های درون توزیع آموزش داده شده است (قسمت آبی) و آموزش قسمت قرمز که وظیفه‌ی تشخیص داده‌هایی که از دسته‌های بدیع آمده اند را دارد
			\cite{Hendrycks2016ABF}.
		}	
		\label{fig:openset1}
	\end{center}
\end{figure}

در کار دیگری با اشاره به اینکه نمی‌توان در تشخیص مجموعه باز، یک دسته ناشناخته اضافه کرد و اگر داده‌ای از دسته‌ی بدیع بود به آن اختصاص داد، یک لایه به نام 
OpenMax
معرفی می‌کند و با اعمال آن در لایه‌ی 
\trans{ماقبل آخر}{Penultimate}
 یک شبکه‌ی عصبی، سعی در تشخیص مجموعه باز دارد
 \cite{Bendale2015TowardsOS}
 . همچنین با بررسی‌هایی که انجام داده، نشان داده است که با اعمال یک 
 \trans{آستانه}{Threshold}
 بر روی خروجی‌های احتمالاتی نمی‌توان به تشخیص مجموعه باز پرداخت چرا که لایه‌ی سافت‌مکس یک ماهیت مجموعه بسته دارد.
 
 \begin{figure}[!htbp]
 	\begin{center}
 		\includegraphics[width=0.9\textwidth]{images/openset2.png}
 		%		\vspace{0.5in}
 		\caption{
 			مصور سازی تابع انرژی روی نمونه‌های واقعی، دستکاری شده و خارج از توزیع به ازای هر دسته. با توجه به این شکل، مشخص است در داده‌هایی که از دسته‌ی بدیع می‌آیند، خطوط سبز پرتکرار است و به این ترتیب می‌توان آن‌ها را تشخیص داد
 			\cite{Bendale2015TowardsOS}.
 		}	
 		\label{fig:openset2}
 	\end{center}
 \end{figure}

در این کار با تحلیل انرژی یا همان تابع امتیاز لایه قبل از سافت‌مکس در شبکه‌های عصبی، نشان داده اند که تمایز واضحی بین داده‌ی واقعی با داده‌های خارج از توزیع هر دسته وجود دارد که نمونه آن در شکل 
\ref{fig:openset2}
آمده است. با استفاده از همین موضوع، برای دسته‌ها میانگین بردارهای فعالسازی لایه‌ی آخر را حساب کرده اند و برای هر دسته نماینده‌ی انرژی آن را بدست آورده اند. سپس برای یک نمونه‌ی جدید با محاسبه‌ی فاصله‌ی بردار فعال‌سازی آن نمونه تا میانگین، راجع به خارج از توزیع بودن داده اظهار نظر می‌کنند.

\subsection{رویکرد ایجابی؛ تنظیم مدل با اضافه کردن توانایی طبقه‌بندی دسته‌های جدید}

در این روش‌ها سعی می‌شود وزن‌های دسته‌ی جدید در مدل را از روی تعداد نمونه‌ی محدودی که برای دسته‌ی بدیع تشخیص داده شده است، حساب کنند. به طور کلی یک راه‌حل متداول 
\trans{حکاکی}{Imprinting}
 وزن‌ها است
 \cite{Qiao2017FewShotIR,Qi2017LowShotLW}
 . 
  \begin{figure}[!htbp]
 	\begin{center}
 		\includegraphics[width=1\textwidth]{images/openset3.png}
 		%		\vspace{0.5in}
 		\caption{
 		مصورسازی به کمک 
 		t-SNE
 		برای مقادیر فعال‌سازی هر دسته قبل از لایه آخر یک شبکه‌ی عصبی تمام متصل در سمت چپ و پارامتر‌های مربوط به هر دسته در لایه‌ی آخر شبکه‌ی تمام متصل در سمت راست. هر نقطه‌ی آبی یک دسته را نشان می‌دهد. نقاط مشخص شده با شکل یکسان و رنگ یکسان مربوط به یک دسته اند. از طرفی اشکال مشابه دسته‌های درشت‌دانه را نشان می‌دهند. به عنوان مثال مثلث‌ها دسته‌های مربوط به پرندگان اند
 			\cite{Qiao2017FewShotIR}.
 		}	
 		\label{fig:openset3}
 	\end{center}
 \end{figure}

شهود این مسئله در شکل 
\ref{fig:openset3} 
مشخص است. در این مصورسازی نشان داده شده که وزن‌های لایه‌ی آخر یک شبکه‌ی عصبی با 
\trans{تعبیه}{Embedding}
یک داده از همان دسته، در جای نزدیک به هم در فضای تعبیه‌ها قرار می‌گیرند. در واقع ارتباط مستقیم بین مقادیر فعال‌سازی و پارامتر‌های همان دسته وجود دارد. این موضوع با عملکرد شبکه‌های عمیق سازگار است. چرا که برای محاسبه‌ی امتیاز یک دسته، مقادیر فعال‌سازی مربوط به داده‌ی فعلی را در وزن‌های مربوط به نورون آن دسته ضرب داخلی می‌کنیم یا به نوعی شباهت کسینوسی می‌گیریم. در نتیجه مقادیر فعال‌سازی به وزن‌های هر دسته‌ای که شبیه‌تر باشد، به آن دسته اختصاص می‌یابد. با در نظر گرفتن این موضوع، می‌توان برای دسته‌های بدیع، با میانگین‌گیری روی مقادیر فعال‌سازی نمونه‌های آن دسته که از یک شبکه‌ی پیش‌اموزش دیده به‌دست می‌آیند، وزن‌های دسته جدید را حکاکی کرد. این موضوع در شکل
\ref{fig:openset4}
نشان داده شده است.

  \begin{figure}[!htbp]
	\begin{center}
		\includegraphics[width=0.9\textwidth]{images/openset4.png}
		%		\vspace{0.5in}
		\caption{
	یک شبکه‌ی عصبی پایه که با دسته‌های درون دامنه آموزش داده شده است و وزن‌های لایه‌ی آخر آن با مستطیل سبز نشان داده شده است. سپس برای داده‌های خارج از توزیع که از دسته‌های بدیع هستند، از این شبکه پیش‌آموزش دیده به عنوان استخراج‌گر تعبیه استفاده ‌می‌شود و این تعبیه بدست آمده که با رنگ نارنجی نشان داده شده است، به عنوان وزن‌های کلاس جدید اضافه و به نوعی حک می‌شود.
			\cite{Qi2017LowShotLW}.
		}	
		\label{fig:openset4}
	\end{center}
\end{figure}

 

\section{یادگیری سلسله‌مراتبی نمونه محدود}


در رویکرد متایادگیری مبتنی بر نماینده، کارهایی در زمینه‌ی استفاده از ساختار سلسله‌مراتبی در داده‌ها مطرح شده است. یک روش تحمیل ساختار سلسله‌مراتب روی یادگیری فضای ویژگی‌ها است
\cite{Garcia2021LeveragingHS}
.
در این روش، از یک مجموعه دادگان که از قبل ساختار سلسله‌مراتبی روی مسئله‌ی هدف دارد، استفاده می‌کنند. ساختار مدل به این صورت است که در ابتدا برای تمامی دسته‌ها در ریزدانه‌ترین حالت ممکن، نماینده‌ها محاسبه می‌شوند. این نماینده‌ها در درخت سلسله‌مراتب پایین‌ترین سطح اند و در واقع برگ‌های درخت را نشان می‌دهند. سپس با استفاده از سلسله‌مراتب از پیش تعیین شده (مانند استفاده از یک مجموعه دادگان سلسله‌مراتبی)، نماینده‌های یک سطح بالاتر در درخت سلسله‌ی مراتب، با میانگین‌گیری روی نماینده‌های زیرمجموعه‌ی هر دسته به‌دست می‌آیند. این طریقه‌ی محاسبه در شکل 
\ref{fig:hier}
آمده است.

  \begin{figure}[!htbp]
	\begin{center}
		\includegraphics[width=1\textwidth]{images/hier.png}
		%		\vspace{0.5in}
		\caption{
			متایادگیری مبتنی بر متر سلسله‌مراتبی؛ در ریزدانه‌ترین حالت، نماینده‌های دسته‌ها محاسبه می‌شود و به صورت بازگشتی، از روی آن‌ها نماینده‌ی دسته‌های سطوح بالاتر یا متانماینده‌ها محاسبه می‌شود. در زمان استنتاج، دسته‌بندی در همه‌ی سطوح انجام می‌شود و خطا برای سطوح ریزدانه و درشت‌دانه محاسبه و وزن‌ها به روز رسانی می‌شوند
			\cite{Garcia2021LeveragingHS}.
		}	
		\label{fig:hier}
	\end{center}
\end{figure}

در واقع این روش یک بسط روی متایادگیری مبتنی بر نماینده است. در پایین‌ترین سطح درخت با در نظر گرفتن دسته‌های مسئله در ریزدانه‌ترین حالت ممکن، با اعمال متایادگیری مبتنی بر متر، نماینده‌ها را بدست می‌آوریم. سپس به ترتیب به سطح‌های بالاتری می‌رویم و هر بار با اعمال روش متایادگیری مبتنی بر متر، نماینده‌هایی از نماینده‌های دسته‌های سطح پایین می‌سازیم که می‌توان به آن‌ها واژه‌ی 
\trans{متانماینده}{Meta Prototype}
را نسبت داد. در معادله‌ی 
\ref{eq:hierarchy}
تابع خطای این مدل آمده است. 

\begin{equation}
	\label{eq:hierarchy}
	\begin{split}
		\mathcal{L}_{hierarchical} = \sum_{h=0}^{H} e^{-\alpha . h} \mathcal{L}_{CE}^{(h)}
	\end{split}
\end{equation}

در این تابع خطا، یک 
\trans{ابرپارامتر}{Hyperparameter}
برای تنظیم مقدار توجه به هر سطح از سلسله‌ی مراتب وجود دارد که با
$\alpha$
نشان داده شده است. از ضریب حاوی این ابرپارامتر برای تنظیم میزان اثردهی تابع خطای 
\trans{بی‌نظمی متقابل }{Cross Entropy}
در سطوح مختلف درخت سلسله‌مراتب استفاده می‌شود. مقدار 
\lr{h}
برابر صفر نشان دهنده‌ی پایین‌ترین سطح در سلسله‌مراتب (برگ‌ها) می‌باشد. بنابراین مقدار 
$\alpha$
اگر مثبت باشد، توجه را روی دسته‌بندی در سطوح با ریزدانگی بیشتر منعطف می‌کند و خطای بیشتری برای آن سطوح در نظر می‌گیرد. اگر مقدار آن برابر صفر باشد، در تمامی سطوح رفتار مشابهی دارد و اگر مقدار منفی داشته باشد، توجه بیشتری به سطوح با ریزدانگی کمتر دارد. به طور کلی اعمال سلسله‌مراتب در متایادگیری مبتنی بر متر سبب می‌شود فضای معنایی که توسط شبکه‌ی نگاشت ویژگی‌ها یاد گرفته می‌شود، این ساختار سلسله‌مراتبی را داشته باشد و بنابراین اگر یک نمونه‌ی جدیدی که در کلاس‌های پایه مسئله وجود ندارد را برای ارزیابی ورودی دهیم، در نزدیکی نماینده‌ای قرار می‌گیرد که متعلق به آن دسته است.

در یک روش دیگر، با تعریف یک خطای مبتنی بر سلسله‌مراتب، به هدایت نماینده‌های هر دسته پرداخته شده است
\cite{Garnot2020LeveragingCH}
.
در این روش، این خطا بر اساس فاصله در سلسله‌مراتب عمل می‌کند. به این صورت که اگر یک داده در یک دسته‌ی اشتباه توسط مدل قرار گیرد، خطای پایه در شبکه‌ی مبتنی بر نماینده در یک ضریب که از ساختار سلسله‌مراتب به‌دست می‌آید، ضرب می‌شود. به طور مثال، با دسته‌بندی اشتباه یک داده، فاصله‌ی آن دسته در یک سلسله‌مراتب از پیش تعیین شده با دسته‌ی صحیح، به صورت کوتاه‌ترین مسیر در درخت سلسله‌مراتب محاسبه می‌شود و این ضریب در خطای پایه شبکه ضرب می‌شود. در نتیجه اگر اشتباه به نحوی باشد که دو دسته‌ی صحیح و اشتباه زیرمجموعه‌ی یک گره والد در فاصله‌ی نزدیک باشند، اثر خطا کمتر است و در نماینده‌های یادگرفته شده، ویژگی سلسله‌مراتبی بودن منعکس می‌شود.

اخیرا در حوزه‌ی یادگیری بدون نمونه، استفاده از مجموعه برچسب‌های سلسله‌مراتبی مورد توجه قرار گرفته‌است
\cite{Novack2023CHiLSZI}
. همانطور که گفته شد، در دسته‌بندی بدون نمونه نیاز به یک مدل 
\trans{تصویری-زبانی}{Vision-Language Model (VLM)}
داریم که مهم‌ترین آن‌ها CLIP است
\cite{clip}
.
در بررسی‌های انجام شده، مشخص شد این مدل‌ها در یک ریزدانگی خاص می‌توانند عملکرد مناسبی داشته‌باشند و اگر برچسب‌های یک مجموعه داده درشت‌دانه باشد، مدل با خطاهای زیادی مواجه می‌شود. به همین دلیل، از مجموعه دادگان سلسله‌مراتبی استفاده کرده‌اند و به جای برچسب یک دسته‌ی درشت‌دانه، برچسب‌های فرزندان ریز‌دانه‌ی آن را به کار برده‌اند. همچنین اگر یک مجموعه‌ دادگان دارای برچسب‌های سلسله‌مراتبی نبود، از یک 
\trans{مدل زبانی بزرگ}{Large Language Model}
استفاده کرده‌اند و به این وسیله برچسب‌های فرزندان ریزدانه را بدست آورده‌اند. با اعمال این فرآیند، دقت‌های دسته‌بندی بدون نمونه افزایش قابل توجهی داشته‌است. نحوه‌ی کار این رویکرد در شکل 
\ref{fig:chils}
مشخص شده‌است. البته 
\trans{باز بودن نسبت به ریزدانگی}{Open-Granular}
هنوز یک چالش باز در این دسته از مدل‌هاست. 

  \begin{figure}[!htbp]
	\begin{center}
		\includegraphics[width=1\textwidth]{images/chils.png}
		%		\vspace{0.5in}
		\caption{
			متایادگیری مبتنی بر متر سلسله‌مراتبی؛ در ریزدانه‌ترین حالت، نماینده‌های دسته‌ها محاسبه می‌شود و به صورت بازگشتی، از روی آن‌ها نماینده‌ی دسته‌های سطوح بالاتر یا متانماینده‌ها محاسبه می‌شود. در زمان استنتاج، دسته‌بندی در همه‌ی سطوح انجام می‌شود و خطا برای سطوح ریزدانه و درشت‌دانه محاسبه و وزن‌ها به روز رسانی می‌شوند
			\cite{Novack2023CHiLSZI}.
		}	
		\label{fig:chils}
	\end{center}
\end{figure}

یک رویکرد دیگر، استفاده از دو دسته مبدا و مقصد و انتقال دانش است
\cite{Guan2020LargeScaleCF,Li2019LargeScaleFL}
.
یعنی در تنظیمات بزرگ مقیاس، یک سری دسته‌ی مبدا داریم که از آن‌ها به تعداد زیاد نمونه در دسترس است. اما در دسته‌های مقصد، با دسته‌های کاملا متفاوت و جدیدی روبرو می‌شویم که از آن‌ها نمونه‌های محدودی در دسترس است. حال می‌خواهیم یک ساز و کار مناسب برای انتقال دانش بین مبدا و مقصد برقرار کنیم. به همین منظور از سلسله‌مراتبی بودن ذاتی مسائل بزرگ مقیاس استفاده می‌کنند و دسته‌های مبدا و مقصد را در قالب یک درخت سلسله‌مراتب و تحت والدهای مشترک به عنوان ابردسته‌ها، تقسیم‌بندی می‌کنند. برای بدست آوردن این درخت سلسله‌مراتب، نیازی به داشتن ارتباط والد و فرزندی بین دسته‌ها نیست و با استفاده از رویکرد‌های داده محور مانند استفاده از بازنمایی زبانی دسته‌ها، مانند آنچه در شکل
\ref{fig:ddhierarchy}
آمده است، می‌‌توان این کار را انجام داد.
   \begin{figure}[!htbp]
 	\begin{center}
 		\includegraphics[width=1\textwidth]{images/ddhierarchy.png}
 		%		\vspace{0.5in}
 		\caption{
 			تشکیل درخت سلسله‌مراتب با رویکرد داده محور؛ از بازنمایی زبانی دسته‌ها که مفاهیم معنایی در خود دارند استفاده می‌کنیم و با یک روش خوشه‌بندی، این درخت را تشکیل می‌دهیم. منظور از Source دسته‌های مبدا و Target دسته‌های مقصد اند و همینطور Superclassها ابردسته‌ها را نشان می‌دهند
 			\cite{Li2019LargeScaleFL}.
 		}	
 		\label{fig:ddhierarchy}
 	\end{center}
 \end{figure}

با داشتن این درخت سلسله‌مراتب، می‌توان یک مدل بازنمایی ساخت که تعبیه‌‌های با
\trans{قابلیت انتقال}{Transferablility}
 بالاتر بین دسته‌های مبدا و مقصد داشته باشند. به این منظور، مشابه شکل
 \ref{fig:knowter}
 از دسته‌های مبدا استفاده می‌کنیم و برای هر سطح در سلسله‌مراتب، یک دسته‌بند اختصاص می‌دهیم. نکته‌ی قابل توجه این است که  دسته‌بند‌های سطوح بالاتر را روی اطلاعات دسته‌بندهای ریزدانه مشروط می‌کنیم و به نوعی این اطلاعات را ورودی می‌دهیم. در نتیجه مدل یاد می‌گیرد اطلاعات ریزدانه را طوری تولید کند تا قابل استفاده در سطوح درشت‌دانه باشد که بین مبدا و مقصد مشترک اند. در نهایت در زمان آزمون و روی دسته‌های مقصد، از بازنمایی مشترک که با رنگ قرمز در شکل 
 \ref{fig:knowter}
 نشان داده شده‌است، استفاده می‌شود. این بازنمایی اطلاعات سطوح درشت‌دانه‌ی مشترک را در خودش دارد و با یک دسته‌بندی مبتنی بر متر مانند نزدیک‌ترین همسایه، دسته‌بندی انجام می‌شود.

   \begin{figure}[!htbp]
	\begin{center}
		\includegraphics[width=1\textwidth]{images/knowter.png}
		%		\vspace{0.5in}
		\caption{
			وارد کردن اطلاعات سلسله‌مراتب در شبکه‌ی بازنمایی مشترک بین مبدا و مقصد که با رنگ قرمز مشخص شده است. با استفاده از داده‌های دسته‌های مبدا که تعداد بالای نمونه دارند، دسته‌بندهای موجود در مستطیل بنفش آموزش داده می‌شوند و با گرادیانی که از آن‌ها به شبکه‌ی بازنمایی مشترک در مستطیل قرمز برمی‌گردد، تعبیه‌هایی با قابلیت انتقال بین مبدا و مقصد تولید می‌شود
			\cite{Li2019LargeScaleFL}.
		}	
		\label{fig:knowter}
	\end{center}
\end{figure}

\section{جمع‌بندی}

این فصل به معرفی کار‌های پیشین مهم و تاثیرگذار از جنبه‌های مختلف یادگیری بزرگ مقیاس، نمونه محدود اختصاص داشت. در ابتدا به معرفی یادگیری نمونه محدود با استفاده از روش‌های سنتی داده‌افزایی در سطح ویژگی و در سطح نمونه پرداختیم. سپس به معرفی رویکرد متایادگیری که یکی از رویکردهای مورد توجه بسیار در دسته‌بندی نمونه محدود است، پرداختیم. در متایادگیری، ساختار کلی روش‌ها بیان شد و سپس الگوریتم‌های مهم این حوزه مانند متایادگیری مبتنی بر بهینه‌سازی و متایادگیری مبتنی بر متر معرفی شد. 

در ادامه به بررسی بدنه‌های بنیادی و مدل‌های پیش‌آموزش دیده برای رسیدن به کارایی داده پرداختیم. سپس به بررسی کارهای پیشین در حوزه‌ی مجموعه باز بودن که یکی از خواص ذاتی مسئله‌ی یادگیری بزرگ مقیاس است پرداختیم. در نهایت به معرفی و بررسی کار‌های مهم در زمینه‌ی یادگیری سلسله‌مراتبی در تنظیمات نمونه‌‌محدود و بزرگ مقیاس پرداختیم.

هر کدام از این سرفصل‌ها به جنبه‌ی خاصی از تنظیمات بزرگ مقیاس، نمونه محدود اختصاص دارند و روی وجهه‌ی خاصی از مسئله تمرکز کرده اند. در فصل بعدی قصد داریم تا روش پیشنهادی خود را معرفی کنیم که بتواند تمامی چالش‌های این مسئله را در بر بگیرد. برای مقایسه‌ی شفاف بین هر کدام از سرفصل‌ها، جدول 
\ref{table:litreview}
در ادامه آمده‌است که بیان می‌کند مقالات بررسی شده در هر سرفصل چه ویژگی‌هایی از مسئله‌ی مدنظرمان را دارند.



\begin{table}[!htbp]
	\caption{سلام}
	\label{table:litreview}
	\setLTR\centering
	\resizebox{\textwidth}{!}{%
		\begin{tabular}{@{}cccccc@{}}
			\toprule
			\rl{توجه به تغییر زیرجمعیت} & \begin{tabular}[c]{@{}c@{}}\rl{استفاده از سلسله‌مراتب ذاتی}\\ \rl{در دسته‌بندی بزرگ مقیاس}\end{tabular} & \begin{tabular}[c]{@{}c@{}}\rl{در نظر گرفتن تغییرات ریزدانه}\\ \rl{(مرز‌های تصمیم‌گیری باریک)} \\ 
				\rl{بین دسته‌های بزرگ مقیاس}
			\end{tabular} & \rl{داده کارآمدی} & \rl{مولد بودن} &                                    \\ \midrule
			
			\fontspec{Segoe UI Symbol}✗                      & \fontspec{Segoe UI Symbol}✗                                                                                             & \fontspec{Segoe UI Symbol}✗                                                                                                                            & \fontspec{Segoe UI Symbol}✓            & \fontspec{Segoe UI Symbol}✗         & یادگیری نمونه محدود با داده افزایی \\
			\fontspec{Segoe UI Symbol}✗                      & \fontspec{Segoe UI Symbol}✗                                                                                             & \fontspec{Segoe UI Symbol}✗                                                                                                                            & \fontspec{Segoe UI Symbol}✓            & \fontspec{Segoe UI Symbol}✗         & 
			متایادگیری مبتنی بر بهینه‌سازی     \\
			\fontspec{Segoe UI Symbol}✗                      & \fontspec{Segoe UI Symbol}✗                                                                                             & \fontspec{Segoe UI Symbol}✗                                                                                                                            & \fontspec{Segoe UI Symbol}✓            & \fontspec{Segoe UI Symbol}✗         & 
			متایادگیری جعبه سیاه               \\
			\fontspec{Segoe UI Symbol}✗                      & \fontspec{Segoe UI Symbol}✗                                                                                             & \fontspec{Segoe UI Symbol}✗                                                                                                                            & \fontspec{Segoe UI Symbol}✓            & \fontspec{Segoe UI Symbol}✓         & 
			متایادگیری مبتنی بر متر            \\
			\fontspec{Segoe UI Symbol}✗                      & \fontspec{Segoe UI Symbol}✗                                                                                             & \fontspec{Segoe UI Symbol}✓                                                                                                                            & \fontspec{Segoe UI Symbol}✓            & \fontspec{Segoe UI Symbol}✗         & 
			استفاده از بدنه‌ی بیانگر           \\
			\fontspec{Segoe UI Symbol}✓                      & \fontspec{Segoe UI Symbol}✗                                                                                             & \fontspec{Segoe UI Symbol}✗                                                                                                                            & \fontspec{Segoe UI Symbol}✗            & \fontspec{Segoe UI Symbol}✓         & 
			مجموعه باز بودن                    \\
			\fontspec{Segoe UI Symbol}✗                      & \fontspec{Segoe UI Symbol}✓                                                                                             & \fontspec{Segoe UI Symbol}✗                                                                                                                            & \fontspec{Segoe UI Symbol}✓            & \fontspec{Segoe UI Symbol}✗         & 
			یادگیری سلسله‌مراتبی نمونه محدود   \\ \bottomrule
		\end{tabular}%
	}
\end{table}





